%% =====================================================================
%%  resumen-scrum.tex
%%  Hoja de una pagina con la estimacion de esfuerzo y costo bajo
%%  metodologia Scrum. Documento autonomo: no usa si3apa7.cls ni entra
%%  en main.tex.
%%
%%  Compilar:  make estimacion-scrum
%%  Salida:    entregables/Estimacion Scrum de esfuerzo y costo.pdf
%% =====================================================================
\documentclass[10pt,letterpaper]{article}

\usepackage[
  letterpaper,
  top=1.6cm, bottom=1.4cm, left=1.9cm, right=1.9cm
]{geometry}
\usepackage{fontspec}
\IfFontExistsTF{Arial}
  {\setmainfont{Arial}\setsansfont{Arial}}
  {\setmainfont{Liberation Sans}\setsansfont{Liberation Sans}}
\usepackage[spanish,es-tabla,es-noshorthands]{babel}
\usepackage{booktabs}
\usepackage{array}
\usepackage[table]{xcolor}
\usepackage{enumitem}
\usepackage{ragged2e}
\usepackage{tikz}
\usetikzlibrary{arrows.meta,positioning,calc}
\usepackage[hidelinks]{hyperref}

\pagestyle{empty}
\setlength{\parindent}{0pt}
\setlength{\parskip}{2pt}
\renewcommand{\arraystretch}{1.03}
\setlist{nosep,leftmargin=1.1em}

\definecolor{azul}{HTML}{1F3A5F}
\definecolor{gris}{HTML}{555555}
\definecolor{linea}{HTML}{BBBBBB}
\definecolor{borde}{HTML}{5A6B80}
\definecolor{relleno}{HTML}{F2F5F9}

% --- Encabezados ------------------------------------------------------
\newcommand{\seccion}[1]{%
  \par\vspace{6pt}%
  {\color{azul}\bfseries\large #1}\par
  \vspace{1pt}%
  {\color{linea}\hrule height 0.6pt}\vspace{3pt}%
}

\newcolumntype{L}[1]{>{\RaggedRight\arraybackslash}p{#1}}
\newcolumntype{R}[1]{>{\RaggedLeft\arraybackslash}p{#1}}
\newcolumntype{C}[1]{>{\centering\arraybackslash}p{#1}}

% --- Cajas del flujo del sprint ----------------------------------------
\newcommand{\caja}[3][]{%
  \node[draw=borde, fill=relleno, rounded corners=2pt, line width=0.5pt,
        inner sep=2.5pt, align=center, font=\scriptsize, #1] (#2) {#3};}

\begin{document}

% =====================================================================
{\color{azul}\bfseries\LARGE Estimación de esfuerzo y costo}\par
\vspace{1pt}
{\large Metodología Scrum \textperiodcentered\ Sistema web de gestión de paquetes
turísticos y reservas \textperiodcentered\ Vizcacha Tours, Potosí}\par
\vspace{2pt}
{\small\color{gris} Sprint 1 del 12 al 25 de noviembre de 2026 \textperiodcentered\
25 de noviembre de 2026 \textperiodcentered\ 1,38 mes-persona \textperiodcentered\
6.875 Bs}\par
\vspace{2pt}
{\color{linea}\hrule height 0.6pt}\vspace{3pt}

% =====================================================================
\seccion{1. El sprint 1 y su compromiso}

{\small
\begin{tabular}{@{}L{0.36\textwidth}L{0.30\textwidth}L{0.28\textwidth}@{}}
\toprule
\rowcolor{azul!8}
{\color{azul}\bfseries Concepto} & {\color{azul}\bfseries Cálculo} &
{\color{azul}\bfseries Valor} \\
\midrule
Duración & 10 días hábiles & del 12 al 25 de noviembre \\
Equipo & 2 desarrolladores (RE y EV) × 4 h por día & 8 h-persona por día \\
Capacidad del sprint & 10 días × 8 h & 80 h-persona \\
Capacidad en puntos & 80 h ÷ 4 h por punto & 20 puntos \\
Compromiso inicial & HU-08, HU-01 y HU-04 & 8 puntos \\
Ampliación a mitad de sprint & HU-02, HU-03, HU-10 y HU-05 & +9 puntos \\
\textbf{Compromiso del sprint} & 8 + 9 & \textbf{17 puntos} \\
Estado al 17 de noviembre & 7 tareas terminadas, 4 historias en curso & 8 pts entregados \\
\bottomrule
\end{tabular}}

\vspace{2pt}
{\small\color{gris} El compromiso ocupa el 85 \% de la capacidad del sprint: hay
margen para imprevistos, pero no para ampliar más.}

% =====================================================================
\seccion{2. El flujo del sprint}

\begin{center}
\begin{tikzpicture}[
  font=\scriptsize,
  flecha/.style={-{Stealth[length=3pt,width=3pt]}, draw=borde,
                 line width=0.5pt, shorten >=2pt, shorten <=2pt},
  rotulo/.style={font=\tiny\itshape, text=gris, inner sep=1pt}]
  \caja{bl}{Backlog del producto\\19 historias \textperiodcentered\ 55 puntos}
  \caja[right=14mm of bl]{sp}{Sprint 1\\7 historias \textperiodcentered\ 17 pts}
  \caja[right=10mm of sp]{ec}{En curso\\4 historias}
  \caja[right=10mm of ec]{re}{Revisión}
  \caja[right=10mm of re]{he}{Hecho\\8 pts \textperiodcentered\ incremento}

  % Frontera del sprint: línea vertical punteada entre el backlog y el sprint.
  \coordinate (frontera) at ($(bl.east)!0.5!(sp.west)$);
  \draw[gris, dashed, line width=0.6pt]
    ($(frontera)+(0,-0.62)$) -- ($(frontera)+(0,0.62)$);

  \draw[flecha] (bl.east) -- node[rotulo, above, xshift=9pt] {compromete} (sp.west);
  \draw[flecha] (sp.east) -- node[rotulo, above] {toma} (ec.west);
  \draw[flecha] (ec.east) -- node[rotulo, above] {termina} (re.west);
  \draw[flecha] (re.east) -- node[rotulo, above] {acepta} (he.west);
  \draw[flecha, dashed] (he.south) -- ++(0,-0.62)
        -- node[rotulo, below] {se planifica el sprint 2} ($(bl.south)+(0,-0.62)$)
        -- (bl.south);
\end{tikzpicture}
\end{center}

\vspace{-4pt}
{\small\color{gris} La línea punteada marca la frontera del sprint: lo que está
a su derecha se entrega el 25 de noviembre como incremento, y lo que queda a
la izquierda vuelve al backlog del producto.}

% =====================================================================
\seccion{3. La estimación del alcance completo}

{\small
\begin{tabular}{@{}L{0.34\textwidth}L{0.32\textwidth}R{0.28\textwidth}@{}}
\toprule
\rowcolor{azul!8}
{\color{azul}\bfseries Magnitud} & {\color{azul}\bfseries Cálculo} &
{\color{azul}\bfseries Valor} \\
\midrule
Esfuerzo unitario & 32 h ÷ 8 puntos entregados & 4,00 h por punto \\
Esfuerzo del alcance & 188 h + 8 pts × 4,00 h & 220 h-persona \\
En mes-persona & 220 h ÷ 160 h & 1,38 mes-persona \\
\textbf{Costo} & 1,38 × 5.000 Bs & \textbf{6.875 Bs} \\
Velocidad por sprint & 80 h ÷ 4 h por punto & 20 pts \\
Sprints del alcance & 55 pts ÷ 20 pts por sprint & 2,75 sprints \\
Cierre estimado & 27,5 días hábiles de trabajo & 07/01/2027 \\
Costo del sprint 1 & 0,50 mes-persona × 5.000 Bs & 2.500 Bs \\
\bottomrule
\end{tabular}}

\vspace{2pt}
{\small\color{gris} La velocidad de 20 puntos es la que corresponde a la
capacidad completa del sprint: se confirma en la revisión del 25 de
noviembre. Si el sprint cierra por debajo, los sprints del alcance suben y
la fecha se retrasa.}

% =====================================================================
\seccion{4. Qué se revisa el 25 de noviembre}

\begin{minipage}[t]{0.49\textwidth}
{\small\textbf{Riesgos de la estimación}\par}
\vspace{2pt}
{\small
\begin{itemize}
\item \textbf{Velocidad sin medir.} 20 puntos es la capacidad, no un
  promedio observado.
\item \textbf{Muestra mínima.} El ritmo se midió en 4 días hábiles.
\item \textbf{Poco margen.} El compromiso ocupa el 85 \% de la capacidad.
\item \textbf{Ampliación a mitad de sprint.} Admisible mientras no
  comprometa el objetivo del sprint.
\end{itemize}}
\end{minipage}\hfill
\begin{minipage}[t]{0.49\textwidth}
{\small\textbf{Métricas de la revisión}\par}
\vspace{2pt}
{\small
\begin{itemize}
\item \textbf{Velocidad real:} puntos cerrados en el sprint.
\item \textbf{Burndown:} puntos restantes frente a los días que quedan.
\item \textbf{Trabajo sin terminar:} lo que se arrastra al sprint 2.
\item \textbf{Objetivo:} si el MVP está completo y es usable.
\end{itemize}}
\end{minipage}

\vspace{4pt}
{\color{linea}\hrule height 0.6pt}\vspace{3pt}
{\footnotesize\color{gris}
Estimación fechada el 17 de noviembre de 2026. Unidades: 1 día-persona = 4 h;
1 mes-persona = 160 h (tiempo completo); tarifa = 5.000 Bs por mes-persona
(31,25 Bs/h); tamaño relativo = 1 punto de historia (escala Fibonacci).
Para recalcularla:
\texttt{Esfuerzo restante = puntos pendientes × (horas-persona gastadas ÷ puntos cerrados)}.
No cubre el mantenimiento posterior al despliegue, ni el costo del hardware,
ni el esfuerzo de documentación y pruebas como tareas aparte.}

\end{document}
